\chapter*{Preface}
Every program manipulates finite objects: bits and formulas, lists to be ordered, networks of roads and wires,
tables of who prefers whom. This book is about the mathematics of those objects, and about a surprise that runs
through it: a few simple ideas, applied with care, settle questions that look hopelessly large. Counting the same
thing in two ways, following an augmenting path, looking at an eigenvector, or bounding the dimension of a vector
space decides problems where trying every case would outlast the universe.

The book starts with logic, the language in which a computer is told what is true, and with orders, the
structures behind sorting, scheduling and termination proofs. It then counts: trees, paths, necklaces and
subspaces, by bijections, by inclusion--exclusion over a lattice, and by generating functions, which turn a
recursive description of a family into an equation and its asymptotic size into the study of a singularity. The
heart of the book is graph theory: Euler tours and de Bruijn sequences, Hall's marriage theorem and the max-flow
min-cut theorem, Kuratowski's characterisation of planar graphs and the five colour theorem, Ramsey's theorem
that complete disorder is impossible, and the spectrum of a graph, which counts its spanning trees and finds its
bottlenecks. The last chapters show three powerful general methods: matroids, the exact setting where greedy
algorithms work; the linear algebra method, which leads to the recent solution of the cap set problem; and
finite geometries, whose designs plan field trials, walk schoolgirls and test software.

The only prerequisite is the toolkit of Tome~0: induction, binomial coefficients, linear algebra and a little
analysis. A few results from other tomes are recalled where they are used, with a precise reference. Every
theorem is proved, and the few that are quoted for their beauty, such as the four colour theorem, are labelled as
remarks with a reference to a complete proof.

Each chapter opens with a question to think about before reading on, and ``Pause and think'' boxes ask you to
guess before a key result. Exercises are graded \textbf{$\star$} warm-up, \textbf{$\star\star$} for a good
afternoon and \textbf{$\star\star\star$} gems; each has a hint and a complete solution that starts with
\emph{How to find it}, because finding the idea matters more than checking the details.

The computer is part of the story. Each chapter has a companion Python notebook that builds the chapter's
algorithms from scratch and tests them against brute force: a satisfiability solver that is then used to find
Ramsey colourings and van der Waerden numbers, generators of random trees that confirm the predicted shapes,
matching and flow algorithms, a spectral clustering of a social network, and a generator of test suites from
finite geometry. Running them is the best way to believe the theorems, and to see where they stop.
